\chapter{Conclusioni}
Il presente lavoro di tesi, ha permesso di affrontare l'intero ciclo
di sviluppo del software. Sono stati curati tutti gli aspetti,
dalla raccolta dei requisiti e dalla progettazione dell'architettura fino alla
realizzazione del Model e della View.
L'obiettivo principale è stato creare un programma
di vendita che potesse sostituire l'attuale sistema. La raccolta dei requisiti è stata fondamentale per comprendere
le necessità degli utenti e le funzionalità richieste, ma
anche per definire cosa non fosse necessario implementare. La presenza dell'autore come 
dipendente in azienda ha rappresentato un vantaggio
significativo per il progetto, permettendo al contempo di sfruttare le conoscenze
pregresse e di interfacciarsi direttamente con gli utenti finali.
Quest'ultimo aspetto ha guidato la scelta della progettazione agile,
che ha permesso di avanzare iterativamente e raccogliere feedback continui dagli utenti.
L'applicativo è stato sviluppato utilizzando il framework
\dotnet{} in tutti gli strati della soluzione. Esso gestisce infatti, mediante \ac{EF Core},
la logica di business e di accesso ai dati e mediante \maui{},
la presentazione grafica.
Tale sistema offre inoltre la possibilità di creare applicazioni cross-platform.
Tra i vantaggi riscontrati nel mantenere un unico framework vi è la possibilità
di condividere il codice tra i diversi strati dell'applicazione,
semplificando lo sviluppo e la manutenzione del software. L'ambiente \dotnet{} ha offerto inoltre una vasta 
gamma di librerie e strumenti che,
attraverso l'uso di \nuget{}, hanno permesso nelle ultime fasi progettuali di
integrare moduli aggiuntivi come la generazione di report. La progettazione del database, attraverso l'utilizzo di \ac{SSIS},
ha consentito di costruire la struttura dati con metodologie professionali,
garantendo la possibilità di apportare modifiche in modo agile.
Nel modulo di interfaccia al database, la scrittura degli Unit Test
ha richiesto in una prima fase un notevole sforzo, ma ha permesso di individuare
e correggere tempestivamente eventuali errori o incongruenze nel codice.
Le modifiche al database sono state costantemente validate dai test, prima
di procedere all'aggiornamento degli strati superiori, garantendo la stabilità e
l'affidabilità del sistema. La realizzazione delle interfacce grafiche è stata
facilitata dall'utilizzo del pattern Factory. Le classi di utilità per i componenti grafici,
hanno permesso di sviluppare agevolmente nuovi ambienti,
mantenendo la coerenza visiva e funzionale dell'applicazione.
Parallelamente, l'adozione del pattern \ac{MVVM} ha consentito di separare
la logica di presentazione dalla logica di business, migliorando la manutenibilità
e facilitando la gestione degli eventi attraverso il data binding.
Allo stato attuale l'applicazione è prossima a un primo rilascio.
Nella versione desktop è già possibile eseguire l'intero ciclo di vendita,
consentendo all'utente finale di:
\begin{enumerate}
    \item accedere al programma tramite autenticazione;
    \item cercare o creare un nuovo cliente;
    \item generare un preventivo personalizzato in base alle esigenze del cliente;
    \item generare eventuali versioni alternative del preventivo;
    \item creare un report in formato PDF del preventivo;
    \item convertire il preventivo confermato in un ordine di vendita;
    \item associare una destinazione di consegna all'ordine di vendita.
\end{enumerate}
Per quanto riguarda la parte amministrativa,
l'applicazione consente, oltre a quanto sopra descritto,
di gestire i prodotti, creare configurazioni, modificare le tabelle con i
dati precompilati destinati agli utenti finali e infine amministrare i permessi
di accesso al programma. Negli sviluppi futuri si prevede di realizzare le interfacce grafiche per
la parte mobile e di implementare un database locale SQLite,
in modo da utilizzare l'applicazione anche in modalità offline.
Ulteriori sviluppi saranno affrontati sulla base dei feedback successivi
ai primi rilasci.